\section{Choosing a future}\label{sec:minds-future}

An engaging conversation offers evidence about a person. It does not reveal how that person will respond to illness, share work, manage disagreement, or care for a child. Someone deciding whether to meet again faces one forecast; someone deciding whether to share a household faces another. The information available, the cost of a mistake, and the ability to revise a decision differ. A useful account of partner assessment must explain these differences before it asks how assessment could influence human evolution.

We organize the problem around future contribution. Capacity to provide resources, knowledge, pleasure, care or protection benefits a particular person through directed investment and compatibility, subject to expected harm. Capacity may be apparent before commitment; reliability requires opportunities to follow through. These components motivate a testable model, without assuming a universal psychology or conscious calculation.

Three arguments follow: shared priorities coexist with differences in emphasis; added observation helps when it contributes independent information; and assessment affects outcomes through particular opportunities and decision rules. Agency matters because assessors may not control marriage, resources, residence or departure. A community can inform a choice while constraining it.

The empirical sequence distinguishes priorities, received ratings, decisions, intentions and outcomes. Each answers a different question; replacing one with another requires validation. Agreement about a presentation does not establish accuracy about later contribution.

\subsection{A global comparison with visible limits}

Exploratory secondary analyses connect international surveys, historical platforms, the Columbia meeting experiment and U.S. GSS family outcomes. These sources concern different populations, rather than a linked longitudinal sample. Their value lies in identifying what each measure establishes and which further observations an inference requires.

``Global'' describes recorded settings, without implying world-population representativeness. Principal international summaries weight eligible country samples equally and retain profile-weighted sensitivities. The estimand remains an average over included samples. Recorded residence, derived residence and survey recruitment location remain distinct; none establishes nationality or ancestry.

The supplied drafts and already inspected platform results informed the analysis plan. The expanded analyses are exploratory, with their specifications and prior visibility recorded in the evidence package. The principal age contrast compares 25--29 with 45--49; country support is held fixed across that comparison. All assigned overview panels remain visible when a result is small, heterogeneous, or contrary to the motivating account. The supplement provides accounting, country atlases, sensitivities, formal derivations, and the complete frozen experiment progression.

Terminology follows the recorded variables. Some sources call their categories sex and others gender. We refer to women and men only for the reported groups actually analyzed; we distinguish the chooser from the person rated. Most available observations concern heterosexual or opposite-recorded-gender relationships. They cannot establish the preferences or experiences of every identity or orientation. Counts of children are variables for studying reproduction, not measures of a person's value. These boundaries define which comparisons the data can support.

\section{Shared priorities, different emphases}\label{sec:minds-priorities}

\subsection{A fixed budget makes trade-offs explicit}

The Kowal project assembled 117,293 responses from 175 countries and included a task allocating 30 points across six long-term partner qualities: health, kindness, physical attractiveness, financial prospects, religion, and an appropriate age \citep{kowal2025,kowaldata}. We reproduce the newer draft's adult heterosexual man/woman cohort with complete, in-range allocations summing to 30, then apply a common country support rule. The retained cohort and country counts are generated from the deposited file, rather than copied from a headline sample size.

\input{minds-kowal-results.tex}

A fixed total matters for interpretation. Additional points for attractiveness necessarily leave fewer for some combination of the remaining options. A difference is an emphasis within that menu, rather than an independent marginal effect of a trait on choosing a partner. The task does not separately measure intelligence, trust, courage, or caregiving skill. Nor does allocating points to an appropriate age specify whether the desired partner is younger or older. Religion's importance and the desired level of religious belief are different questions. These restrictions prevent a seemingly simple ranking from becoming a complete theory of mate value.

Figure~\ref{fig:minds-preferences} retains every budget item alongside the Australian ratings. Country- and respondent-weighted means answer different questions; neither weighting repairs convenience-sample recruitment.

\mwcfigure{minds-preferences.pdf}{fig:minds-preferences}{\textbf{Stated priorities on their original scales.} A: all six Kowal budget items, with equal country weights and respondent-weighted means for comparison. Each person's budget totals 30. B: all nine Whyte importance items in the deposited cohort, retaining openness and competing priorities. C: selected Whyte age profiles; complete items, four age bands and paired contrasts are in the supplement. Women and men identify the respondents. Error bars describe the specified sampling calculation and do not establish national representativeness.}{preferences}

\subsection{Variation shared by women and men}

The average woman--man gap omits variation shared by both groups. Country-sample means for religion and kindness covary across genders, compatibly with differences in norms, recruitment and composition. Complete estimates retain magnitudes, uncertainty, reversals and sample support; they do not identify one cultural cause.

\input{minds-kowal-context.tex}

The source used local collaborators, online and social recruitment, outsourcing platforms and some paper questionnaires \citep{kowal2025}. Translation, mode and composition can affect country differences. Age standardization addresses one compositional dimension within supported samples; a country label does not identify a stable psychological type.

Walter and colleagues' 45-country project provides a second international source on five desired qualities measured from 1 to 7 \citep{walter2020}. We reconcile code 0 with women and code 1 with men using their exact matches to the primary article's final sample composition: 14,399 participants, including 7,909 women. This is an inference from documented counts, independent of preference outcomes; the original literal code key remains unavailable. \input{minds-walter-results.tex} Agreement between surveys is informative after their questions and coding are aligned, rather than because both contain an attractiveness variable.

\mwcfigure{minds-countries.pdf}{fig:minds-countries}{\textbf{Country variation is part of the result.} Named supported country samples, their composition and selected preference and age contrasts are shown with a fixed contrast orientation. Country-level effects remain visible rather than being replaced by counts of positive signs. The atlases include every supported country and sample size. Estimation intervals and between-country dispersion answer different questions.}{countries}

\subsection{Trust, intelligence, appearance, and competing items}

Whyte and colleagues asked Australian respondents to rate nine attributes spanning aesthetics, resources, and personality \citep{whyte2021}. The publicly deposited records and the article's headline sample are not interchangeable: the source-flow audit states which records are present and eligible, and does not invent an explanation for any unreconciled difference. The reanalysis preserves all nine items, four age bands, the documented heterosexual subset, and within-person centered responses. Paired trust-minus-intelligence and trust-minus-attractiveness comparisons use respondents who answered both items.

\input{minds-whyte-results.tex}

Importance scales allow several traits to receive high scores. They therefore complement rather than replicate a forced budget. Trust can be highly valued while physical attraction remains influential in an encounter. Within-person centering addresses general differences in scale use, but changes the estimand to relative emphasis within each person's nine responses. It does not isolate an objective value of trust or transform a survey report into observed dependable behavior. Age profiles also compare respondents of different ages, rather than following an individual's priorities through life.

The item wording also differs: the intelligence question refers to deciding to be sexually intimate, whereas several other items concern sexual attractiveness. Paired numerical contrasts preserve those questions; they do not equate their psychological targets.

The broader lesson is a conjunction. Shared priorities can coexist with recurring average differences; both can vary with the decision and the measurement. A finding that women allocate more to financial prospects in a particular task does not imply indifference to appearance, just as men's emphasis on appearance does not imply indifference to kindness. Complete distributions and effect sizes are more informative than a categorical description of what either gender wants.

\section{How profiles are rated}\label{sec:minds-ratings}

Historical HON records permit a different question: how do received appearance ratings vary with the profile holder's age, recorded gender, country, and voting exposure? We scan the complete canonical metadata population before applying eligibility. The primary cohort requires ages 18--79, valid parsed vote and like counts, unconflicted metadata and labels, a recognized recorded gender, and at least 20 votes. A finite supplied beauty score is additionally required for score and percentile summaries. The source register distinguishes raw source rows, canonical profiles, and each analysis denominator.

\input{minds-hon-results.tex}

\subsection{An age curve needs a reference population}

Raw beauty scores and like rates preserve their original units. Percentiles provide another view of where profiles fall in a specified distribution. Our principal percentile is the tie-adjusted midrank within country and recorded gender, using one adult reference spanning the eligible ages. That reference stays fixed across age. Re-ranking separately at each age would force every age group to have approximately the same mean percentile and erase the comparison of interest. The original global within-gender and pooled-gender references remain sensitivities.

For principal age comparisons we hold the country set fixed: at least 100 eligible profiles of each recorded gender in each of ages 18--29, 30--44, and 45--64. Normalization of country strings is versioned and unresolved labels remain explicit. Country-specific natural cubic spline fits cover ages 18--64 with interior knots at 25, 35, 45, and 55. The contrast averages predictions over exact ages 25--29 and 45--49, then weights supported countries equally. Observed broad-band means provide an immediate check against the fitted description. Five-year observed summaries through age 79 remain in the supplement.

\mwcfigure{minds-beauty.pdf}{fig:minds-beauty}{\textbf{Received appearance ratings across age.} A--B: scores and fixed country-and-gender-referenced percentiles, with fitted curves and broad-band observations for \MwcHonCountries{} countries. Shading gives 95\% country-bootstrap intervals (2,000 resamples); Figure~\ref{fig:minds-countries} shows country dispersion. C: mean country score quantiles in \MwcHonDistributionCountries{} countries with at least 100 profiles per gender in both narrow bands; thin lines span the 10th--90th quantiles, thick bars the 25th--75th, and points are medians. D: upper-tail shares on the principal set. These compare historical profiles, not within-person aging.}{beauty}

These curves describe participating profiles and recorded audiences. Age covaries with presentation, participation, exposure and selection, while rating dates and voter identities are insufficiently observed. The gradient therefore cannot isolate biological aging, desired partner age or changes in an existing partner's valuation.

Raw distributions, quantiles and upper-tail shares retain information hidden by mean percentiles. Distributions overlap, and country contrasts can reverse. The country bootstrap resamples included countries; it cannot account for missing populations.

\subsection{Attention changes the observation process}

Votes are not an incidental precision variable. A profile receiving many votes has had more opportunities to be seen, while an unobserved selection process determines both who is presented and who votes. We distinguish the mean of profile-specific like rates from total likes divided by total votes. The first weights people equally; the second weights votes equally and gives highly exposed profiles greater influence. Differences between them reveal how attention enters the summary, without identifying why attention was allocated.

\mwcfigure{minds-attention.pdf}{fig:minds-attention}{\textbf{Attention and saved-image coverage.} A--B: equal-country means of country medians and like rates on the principal \MwcHonCountries{}-country set. C: pooled inventory coverage within U.S. and other recorded groups. D: source-wide unadjusted median exposure by reported photo count; adjusted associations are supplementary. Gender refers to profile holders. Metadata eligibility, inventoried files, decoded images and modeling participation remain distinct. Vote thresholds of 1, 20 and 100 are retained as sensitivities.}{attention}

Saved images are geographically concentrated, so global metadata coverage cannot validate a model trained on a predominantly U.S. subset. The audit distinguishes eligible metadata, inventoried images, decoded files, modeling inclusion, completed processing and finalized duplicate components.

Portfolio size can reflect engagement, tenure, presentation or access. Adjustment for age, country, gender and voting history remains associational. Reported photo counts and inventoried files are retained separately because archival differences need not represent participant choices.

HON biography, education/work, language, interest and purpose fields lack overlap with valid labels at the principal vote threshold in the inspected exports. Their schema presence therefore cannot support multimodal rating analysis. Global appearance findings use metadata and supplied ratings; image prediction remains a separate supplementary exercise.

\section{From ratings to decisions}\label{sec:minds-decisions}

In the Columbia speed-dating experiment, participants rated partners and decided whether to request further contact after short meetings \citep{fisman2006}. Its 8,378 directed meeting rows describe 551 people in 21 sessions. Rows are neither independent people nor necessarily complete reciprocal pairs. The accounting distinguishes usable directed decisions, participants, sessions, and matched reciprocal records. The experiment improves on a profile archive for studying a decision after interaction, while retaining a narrow population and an artificial encounter structure.

We consolidate the behavioral analysis around two questions. Does a chooser say yes more often to partners whom other session participants rate more attractive? And does that association remain after accounting for the chooser's perceptions of other qualities? For the first, each partner receives a leave-one-chooser-out mean appearance rating from at least three other session participants. Partners are ranked within each chooser's available dates, with ties assigned midpoint ranks. Five equal-width rank bands summarize yes rates. A band describes relative standing in that chooser's actual opportunity set, rather than a population attractiveness percentile.

\input{minds-speed-results.tex}

\mwcfigure{minds-choice.pdf}{fig:minds-choice}{\textbf{Appearance and the decision to meet again.} A: yes rates across within-chooser appearance-rank bands, using ratings from other session participants. B: adjusted appearance associations from chooser fixed-effect linear probability models using own ratings and leave-one-chooser-out panels. Intervals resample sessions with shared draws for gender contrasts. C: agreement about appearance and perceived intelligence from the supplied v9 analysis, explicitly labeled as agreement. Complete coefficients, raw-gradient uncertainty, session deletion and common-complete-case checks appear in the supplement.}{choice}

The adjusted models use a fixed effect for each chooser and six perceived qualities: attractiveness, fun, intelligence, ambition, sincerity, and shared interests. Within-chooser estimation asks which partners a given participant is more likely to accept among the dates observed. Own-rating and panel-rating models differ in what they condition on, their complete cases, and their error structure. Own ratings may be influenced by a general favorable reaction to the same partner. Panel ratings reduce that particular shared source but remain impressions formed in the same social setting. Neither model randomly changes a partner's appearance or identifies the causal effect of a physical trait.

Whole-session resampling retains the dependence created by a common event and uses the same bootstrap draws for women--men contrasts. Leave-one-session-out estimates reveal whether a small number of sessions drives an adjusted contrast. A common-complete-case sensitivity reduces the possibility that an own-versus-panel difference merely reflects different rows. These checks cannot repair the experiment's selectivity: Columbia graduate students, short opposite-gender encounters, and requests for contact are a specific setting and outcome. A yes does not establish a later meeting, a relationship, or long-term investment.

\subsection{Agreement, desirability, and compatibility}

The older v9 analysis found stronger observer agreement about appearance than perceived intelligence in these encounters. We retain that result as agreement. If several observers share a stereotype, favor a fluent presentation, or all meet someone on an unusually good evening, their ratings may converge without becoming more accurate about an independently measured ability. Calling that convergence ``reading a mind'' would require a criterion that this analysis lacks.

Agreement about targets also differs from chooser selectivity and pair-specific attraction. A generally desired person receives more favorable evaluations across observers. A generally selective chooser says yes less often across targets. A particular dyad may nevertheless depart from both averages. Distinguishing these components matters because a model can predict average desirability while doing poorly at forecasting which two people will form a mutually satisfying relationship. Prospective tests of compatibility need repeated partners or interactions and an outcome defined for the pair.

Eastwick and colleagues' international preference-matching project offers a complementary comparison of stated ideals and romantic evaluations \citep{matching2025}. We retain its full published 35-attribute profile, the original ideal-rating scale, and the published standardized evaluation associations. Those two columns cannot be compared as if their numerical units were interchangeable. The study reports stated gender differences for attractiveness and earning potential that do not reproduce as the corresponding gender differences in evaluation associations. This makes it a relevant boundary for generalizing the Columbia adjusted contrast, rather than an inconvenient result to omit.

The combination supports no single ranking of what matters in every relationship. A trait can be highly valued, poorly recognized, influential in an early choice, and weakly predictive of later relationship quality. Conversely, an initially unremarkable quality may matter greatly when tested by shared work or stress. The next question is which future people describe themselves as seeking, because the anticipated relationship changes which contributions and risks are relevant.

\section{Commitment, children, and timing}\label{sec:minds-goals}

POV contains \MwcPovN{} account-unique historical records with separate fields for relationship intention, the kind of relationship sought, reported parenthood, and desire for children. A field called \texttt{intent} and another called \texttt{for} are separate questions. Their categories can express different aspects of a goal and must not be pooled into an invented commitment scale. The primary analysis retains ages 18--79 and profiles seeking the other recorded gender, with all-seeking and single-profile sensitivities. ``All my kids are over 18'' indicates existing parenthood, rather than childlessness.

Country provenance requires reconstruction because the stored country column is blank. We use a versioned geographic dictionary and hash-pinned Natural Earth boundaries to distinguish unambiguous state or province assignments from coordinate-supported assignments, agreement, conflict, and unresolved cases. Ambiguous province names and contradictory evidence remain unresolved in principal geographic comparisons. Stored coordinates and account identifiers stay at the data source; only aggregates are exported. Derived country is a documented geographic inference about stored residence, rather than a reported nationality.

\input{minds-pov-results.tex}

\subsection{Two intentions can have different age profiles}

Figure~\ref{fig:minds-goals} aligns commitment and childbearing responses within defined cohorts. The age comparison holds country support fixed, with at least 100 profiles in every gender-by-age cell required for the contrast. The denominator is part of the result: reported childlessness, known parenthood, and all eligible profiles describe different populations. We retain affirmative, negative, undecided, and missing responses. An answered-only curve may change because the response process changes, so the all-profile curve remains visible.

\mwcfigure{minds-goals.pdf}{fig:minds-goals}{\textbf{Relationship and childbearing intentions measure different goals.} Aligned POV cohorts show the three documented English \texttt{intent} commitment labels, the separate \texttt{for} responses, desire for children and reported parenthood by age and profile gender. Country-balanced contrasts use the fixed validated geographic support set. Intention percentages retain all responses and nonresponse; parenthood uses known reports. Coordinate-supported checks and complete language-specific response distributions are supplementary. These are cross-sectional historical reports.}{goals}

A lower share wanting additional children in midlife does not by itself indicate less interest in a committed partner. Existing children, biological timing, prior unions, health, and desired family size can change one intention without changing the other. The pooled childbearing reversal in the earlier analysis is retained as a reproducibility baseline. Its interpretation now depends on whether the pattern persists across the validated country set and under exact-age standardization, the coordinate-supported subset, and single-profile and all-seeking sensitivities. A pooled reversal is not a substitute for these comparisons.

The two age bands compare people observed at different ages and historical circumstances. They cannot show that an individual changed their mind or postponed a birth. Platform membership also selects on current availability, relationship history, and willingness to describe oneself publicly. Prior relationships can influence both the probability of being in the archive and the answers given there. A country-balanced summary describes the supported profile samples; it does not recover all adults who formed relationships without joining the platform.

Geographic recovery may select profiles by age, gender, parenthood or intentions; the audit reports those differences. Coordinate--label agreement strengthens provenance without verifying current residence. Boundary and territory conventions are documented.

\subsection{Ambition, education, and relationship history}

Rich self-reports allow a broader descriptive question: how do intention and parenthood patterns vary with stated ambition and education? We specify the outcome, denominator, comparison ages, and adjustment set before fitting. The analysis presents unadjusted associations beside standardized comparisons accounting for age, derived country, marital status, and existing parenthood where appropriate. Existing parenthood is itself the outcome in one analysis and therefore cannot also be used as its adjustment variable. Each model's actual adjustment set is recorded explicitly.

Ambition is a self-description, not a measured cognitive ability, work history, or randomized exposure. Education may precede family formation, follow it, or share causes with it. An association of greater ambition with fewer existing children and greater future childbearing desire would be consistent with several processes, including different ages and constraints. It would not establish reproductive postponement without dates and longitudinal follow-up. The same restraint applies when an adjusted pattern weakens, disappears, or changes direction: those are substantive results about composition, not failed opportunities to tell the preferred story.

The field \texttt{relationshipage\_id} is treated as a categorical report of longest prior relationship duration only to the extent that the question wording is verified. Open-ended categories remain open-ended; ``over ten years'' cannot be replaced by an exact duration. Heights are parsed from explicit centimetre entries with malformed and implausible values audited. Body descriptions, religion, profession, and other self-reported categories remain supplementary comparisons with adequate support. A profession category is not converted into a skill or resource ranking without a documented crosswalk.

These analyses address what future respondents describe as desirable and feasible. They do not observe whether partners agree on that future, whether a stated goal is credible, or whether the profile holder can realize it. Those questions require mutual assessment, later behavior, and the opportunities governing family formation. The transition from intention to outcome is therefore an empirical problem rather than an assumption built into the narrative.

\section{From assessment to outcomes}\label{sec:minds-outcomes}

\subsection{Partnering and parenthood are separate episodes}

Partner formation, becoming a parent, and the number of children among parents are distinct events. A characteristic associated positively with one can be associated negatively with another. An intention to have children may not be realized; an unplanned birth may occur without a stated intention. Contraception, infertility, separation, mortality, resources, and public institutions all intervene. A description of contemporary reproduction therefore needs episode-specific quantities before it can inform an evolutionary interpretation.

The international survey's self-rated qualities and current partnership provide a cross-sectional comparison, rather than a record of how a relationship began. Analyses within countries adjust for the recorded covariates but cannot establish that a self-rating caused partnership. A partner may change a respondent's self-assessment; age and opportunity may affect both. These associations are useful as a check on the assumption that a stated preference directly determines realized partnering. They are not estimates of genetic selection or of an unobserved population's marriage process.

We retain the GSS as a named U.S. case study \citep{gss2024}. The supplied v9 analysis used adults aged 45--79 born in 1900--1979, with known child count and positive survey weight, standardizing traits within recorded sex and birth decade. Personal-income comparisons restricted ages further to reduce retirement-related incomparability. Years of schooling, a short vocabulary score, occupational prestige, and personal income are distinct measured traits. Schooling is not a direct measure of inherited intelligence, and current income does not summarize lifetime resources.

\input{minds-gss-results.tex}

The original weighted decomposition is reproduced before a new design-aware income sensitivity. The latter uses released strata and PSU variables from the hash-pinned cumulative source, resampling PSUs within survey-year-by-stratum groups before restricting to the analysis domain. Its supported subset, singleton-stratum exclusion, 2,000 shared resamples and alternative weight calculation are reported separately. Treating survey years independently remains an assumption. Top-coded child counts, missing traits and nonresponse also limit interpretation. The archived respondent-bootstrap intervals retain their original description and are not relabeled as survey-design intervals.

\subsection{An accounting identity clarifies opposed effects}

Let $c$ be reported child count, $P=1$ indicate $c>0$, and $z$ be a standardized measured trait in a defined weighted population. With positive mean child count, the phenotypic selection differential can be written
\begin{equation}
S=\Cov\!\left(z,\frac{c}{\E[c]}\right)
=\underbrace{\E[z\mid P=1]-\E[z]}_{\text{parenthood component}}
+\underbrace{\frac{\Cov(z,c\mid P=1)}{\E[c\mid P=1]}}_{\text{family-size component}}.
\label{eq:minds-decomposition}
\end{equation}
This is an identity, not a causal partition. The first component describes how parents differ from all adults on the measured trait. It combines chosen childlessness, infertility, timing, mortality and other processes; it must not be interpreted as exclusion from mating. The second describes association with family size among parents. Weighted means and covariances must use the same eligible population for the equality to hold.

\mwcfigure{minds-outcomes.pdf}{fig:minds-outcomes}{\textbf{Outcomes and evolutionary constraints.} A: GSS income differentials separated into parenthood and family-size components. Filled markers and intervals use the newly reproduced, design-supported domain and stratified PSU bootstrap; hollow markers retain the full archived baseline. Their sums follow an accounting identity. B: further empirical links needed between assessment and inherited anatomical change. The associations in A do not identify the pathway in B.}{outcomes}

Opposed episode associations matter even when their sum is small. If a measured trait is positively associated with becoming a parent and negatively associated with family size among parents, reporting only the first would exaggerate the net association; reporting only the total would conceal the mechanism worth investigating. Marriage and parenthood also become less tightly coupled across cohorts. A statement about marital access cannot be transported to a setting where more births occur outside marriage without checking that link.

The same logic limits a one-direction account of modernity. Education, income, partnership, and births can relate differently by gender, cohort, and institution. Register studies in Nordic populations provide relevant comparisons of life-course income and completed fertility \citep{kolk2023,jalovaara2019}. They do not make a U.S. cross-sectional income measure causal or universal. A defensible international comparison would align birth cohorts, ages, child-count definitions, survival, and survey design rather than treating any association with parenthood as interchangeable evidence.

\subsection{What these outcomes contribute to the argument}

Assessment's consequences cannot be read directly from attraction. A valued quality may improve relationship access while accompanying fewer births through timing or competing opportunities; the converse can occur. Cultural influence, cooperation and reproduction require separate measures and may respond differently to institutions.

Women contribute knowledge, earnings, relationships and care directly to household outcomes, as men contribute caregiving and directed effort beyond obtaining a partner. Which delivery matters is an empirical question about particular households, without assigning fixed sexes to resources or appearance.

\section{Learning about dependable contribution}\label{sec:minds-learning}

\subsection{What information is valuable?}

A minimal relationship forecast separates potential benefit $B_j$, the fraction or probability $p_{ij}$ expected to reach chooser $i$, compatibility $C_{ij}$, and expected harm $H_{ij}$:
\begin{equation}
U_{ij}=\E[p_{ij}B_j+C_{ij}-H_{ij}\mid\mathcal I_i].
\label{eq:minds-utility}
\end{equation}
Here $\mathcal I_i$ is the evidence available when the decision is made. Benefits can include affection, understanding, care, practical work, or shared pleasure. The representation is deliberately modest: a model for prospective tests, not a fitted decision law or a claim that reproductive output is people's conscious objective. Different outcomes may require different functions and time horizons.

The model makes one useful distinction exact. Expected contribution is generally
\begin{equation}
\E[pB\mid\mathcal I]=\E[p\mid\mathcal I]\E[B\mid\mathcal I]+\Cov(p,B\mid\mathcal I).
\label{eq:minds-joint}
\end{equation}
A high capacity estimate and a high willingness estimate do not determine their joint contribution when capacity and delivery covary. Someone's most capable behavior may be directed elsewhere. A partner with less spectacular skill may reliably apply it to the relationship. Prospective measurement needs to ask what is delivered, to whom, under which circumstances, and at what cost. A general reputation for competence is only one input.

Institutions can change the same forecast. If a fraction $e$ of an output is reliably delivered independently of willingness, the delivered share becomes $e+(1-e)p$. Holding the model's other terms fixed, willingness then has marginal value $(1-e)B$ while capacity has marginal value $e+(1-e)p$. This comparative static explains why verification needs can change with institutions. It does not show that enforcement eliminates generosity or that coercive arrangements maximize anyone's welfare. Consent, distributional effects, and the costs of enforcement are separate outcomes.

\subsection{Observers and contexts address different errors}

Suppose observer $j$ evaluates target person $i$ in context $k$, and the target is an independently defined, stable contribution $\theta_i$. An illustrative measurement model is
\begin{equation}
Y_{ijk}=\theta_i+h_i+b_{ij}+c_{ik}+\epsilon_{ijk}.
\label{eq:minds-measurement}
\end{equation}
The terms describe a shared persistent halo $h_i$, observer--target distortion $b_{ij}$, context--target disturbance $c_{ik}$, and residual error $\epsilon_{ijk}$. Assume centered, mutually independent components with common finite variances and independent draws across observers and contexts. Normality is unnecessary for the following variance identity.

\begin{proposition}[Observer and context averaging]
For a balanced average over $m$ observers and $k$ contexts under Eq.~\eqref{eq:minds-measurement}, squared correlation with the independent criterion is
\begin{equation}
R^2_{m,k}=\frac{\sigma_\theta^2}{\sigma_\theta^2+\sigma_h^2+\sigma_b^2/m+\sigma_c^2/k+\sigma_\epsilon^2/(mk)}.
\label{eq:minds-validity}
\end{equation}
\end{proposition}

Additional observers reduce observer-specific distortion but cannot remove the shared context of the same performance. Additional contexts reduce context-specific disturbance but cannot remove one observer's persistent distortion. Neither removes shared halo. Two panels can agree partly because they share the same halo; their agreement is not the criterion correlation in Eq.~\eqref{eq:minds-validity}. A photograph rated thousands of times supplies many judgments of one presentation, rather than thousands of independent tests of dependable contribution.

\mwcfigure{minds-assessment.pdf}{fig:minds-assessment}{\textbf{Information, contribution, and opportunity.} A: the reveal, produce and allocate routes have separate observable outcomes and rival explanations. B: illustrative criterion-validity curves under Eq.~\eqref{eq:minds-validity}, varying observers and contexts. The nuisance variances are declared illustration parameters, not estimates from HON or speed dating. C: a prospective design holds observation cost fixed and tests later independent contribution. Shared halo remains a ceiling under the stated model.}{assessment}

The divisors require consequential assumptions. Observers may share training or communicate; contexts may repeat the same presentation; people may change over time. For exchangeable correlated errors, variance reduction slows toward a nonzero common component. Unequal exposure and selective observation also require explicit weighting or a fuller model. Crossed observer and target designs are needed to estimate variation at the relevant levels \citep{judd}. None of our platform analyses estimates the variance components in Eq.~\eqref{eq:minds-validity}.

Context can also be part of the signal. A person dependable in one environment and unreliable in another may be a good match for a particular relationship rather than an error-prone example of a fixed trait. The criterion determines what should be averaged away. Forecasting general performance across settings differs from forecasting one pair's contribution under anticipated constraints. Longer observation helps only if it samples relevant behavior and adds information beyond familiarity, endorsement, and the existing reputation.

\subsection{Revealing, producing, and allocating}

A public activity can reveal qualities, produce benefits, and allocate opportunities. These routes can coincide without being identical. Skilled food acquisition may inform observers, provision others, and alter who receives support or invitations. Storytelling may display memory or social understanding, coordinate behavior, and make someone a preferred social partner. To identify the reveal route, new evidence must improve prediction against an independent criterion. To identify production, the activity must change a recipient or group outcome. To identify allocation, it must change access to a measurable opportunity.

Rivals differ by route: familiarity or endorsement can increase liking without prediction accuracy; resources can produce benefits without assessment; authority can allocate access without skill. Esteem alone cannot distinguish these routes. Measuring outcomes and temporal order makes them testable.

Women's knowledge and influence supply substantive examples. Among Tsimane, parental ethnobotanical knowledge was associated prospectively with child-health indicators, with important maternal associations \citep{mcdade2007}. Another study linked maternal political influence to child growth and illness measures in a particular community after measured covariate adjustment \citep{alami2020}. Neither establishes a universal causal effect or proves a mate-choice pathway. They show why useful knowledge, social ties, and household decision authority belong in an account of contribution. Mothers are participants in production and allocation, not merely objects of assessment.

Cooperative care also broadens the relevant unit. Relatives and other helpers can contribute information and practical support, with effects depending on the relationship and setting \citep{sear2008}. Partner assessment occurs in these networks. A prospective spouse may be evaluated by relatives whose interests overlap only partly with the prospective partners'. The ability to refuse, negotiate, or leave influences whose forecast has consequences. Reliable information cannot make an imposed decision consensual.

\section{Culture and cognitive evolution}\label{sec:minds-evolution}

\subsection{Evidence for each link}

The modern results motivate an assessment hypothesis; they do not estimate its ancestral contribution. The proposed historical pathway requires several links: behavior must contain information about consequential contribution; observers must learn from it; decision-makers must use the information; those decisions must affect cultural transmission or reproductive outcomes; any inherited response must reflect transmissible variation; and anatomical change must be developmentally and energetically feasible. Failure at one link can interrupt the pathway while leaving all earlier associations intact.

Cultural transmission has a comparatively direct route. An admired or useful person can receive more attention, pupils, support, and opportunities to demonstrate a practice. Learners can adopt a technique within a lifetime, without a change in allele frequencies. Prestige accounts distinguish freely conferred deference from force and institutional authority \citep{henrich}. The distinction matters empirically: a high-status person may receive attention because observers judge them informative, because others already defer to them, or because refusal is costly. These mechanisms predict different responses to independent information and changes in authority.

Field studies support parts of the framework without completing the historical chain. In Agta research, skilled storytelling was associated with preferred camp-mate choices and other social outcomes \citep{story}. Such evidence connects a recognized activity with social preference; it does not establish which cognitive component observers detected or whether the preference caused the reported reproductive association. Foraging skill develops over a long period and varies across individuals and settings \citep{koster}. Public success combines knowledge, practice, equipment, health, cooperation, and luck. It is not an isolated readout of inherited general intelligence.

Marriage decision authority is itself variable. The supplied cross-cultural codes distinguish individual choice, family approval, parental arrangement, consultation, consent, and divorce rights \citep{kirby2016,broude1983,frayser1985,whyte1978}. Missingness is substantial and the ethnographic units share histories and neighbors. We retain counts and the original code definitions, with no assumption that every society supplies an independent trial. A daughter being consulted is not equivalent to her consent being necessary; permission to object is not equivalent to an effective ability to refuse. These distinctions limit any simple claim that mutual mate choice universally linked public competence to reproduction.

\subsection{Breadth can favor general contribution under explicit conditions}

The strongest formal idea retained from v9 concerns what a score rewards. Let performance in domain $\ell$ be $X_\ell=a g+s_\ell$, where $g$ is a common capacity and $s_\ell$ a domain-specific component. For $K$ equally weighted domains with independent specialist components,
\begin{equation}
\bar X=a g+\frac{1}{K}\sum_{\ell=1}^{K}s_\ell,\qquad
\Var(\bar X)=a^2\sigma_g^2+\frac{\sigma_s^2}{K}.
\label{eq:minds-breadth}
\end{equation}
Broader comparison can increase the share of score variation attributable to the common component. This is an algebraic property of the stipulated loading and covariance structure, not a demonstration that ancestral activities had those properties or that all people benefited from being generalists. Correlated specialties, unequal weights, division of labor, and choosing one's best task change the result. A community that rewards a single useful specialty may rationally allocate opportunities differently.

Directional selection requires covariance between an inherited component and relative reproductive success, not merely an accurate phenotypic forecast. In quantitative-genetic notation, a response in a vector of traits requires an additive genetic covariance structure as well as selection gradients \citep{lande1983,hazel1943}. Cultural advantage can be transmitted through teaching, resources, institutions, or family circumstances rather than alleles. A phenotypic association between income and children does not identify the genetic covariance that would be needed to predict an inherited response in cognition.

Accumulating small selection differentials in a model cannot identify an ancestral process from uncertain modern correlations. Mortality, fertility regulation, migration, institutions and genetic constraints can alter responses. The supplement retains conditional models and checks; their parameters are not estimates from these surveys or profiles.

\subsection{A larger brain requires support}

Demand for useful information processing is only one side of an anatomical explanation. A developing brain requires energy, time, care, and a viable developmental pathway. Studies of brain metabolism document substantial childhood demands \citep{kuzawa2014}. Models of human life history connect learning and prolonged dependency to intergenerational transfers and foraging returns \citep{kaplan2000}. These constraints mean that a preference for competence cannot itself supply the resources needed to support a larger brain. Cooperative production and care may change feasibility even when assessment is unchanged.

\subsection{Environmental context constrains historical attribution}

The fossil measurements describe cranial volume, rather than a direct measure of cognition or partner assessment. Gingerich's consensus workbook supplies \MwcFossilN{} observations and identifies approximately \MwcStasisOlder--\MwcStasisYounger{} million years ago (Ma) as a published interval of relative stasis \citep{gingerich2022}. Figure~\ref{fig:minds-environment} shows observations and 100-kyr geometric means without imposing a flat trajectory. Apparent phases depend on temporal scale, sampling and model specification. Du and colleagues distinguish whole-clade patterns from processes within lineages; the two need not support the same evolutionary description \citep{du2018}. Ecological and demographic conditions therefore constrain attribution of anatomical change to any proposed social mechanism.

The Middle Pleistocene Transition, conventionally approximately \MwcMptOlder--\MwcMptYounger{} Ma, involved a shift toward longer, stronger glacial cycles rather than one precisely dated catastrophe \citep{clark2006}. The published LR04 benthic oxygen-isotope stack provides broader chronological context: its \MwcClimateRecords{} marine records combine information about global ice volume and deep-ocean temperature, not direct air temperature \citep{lisiecki2005}. Marine and terrestrial proxies document pronounced North Atlantic and southwest European cooling during MIS 34, approximately \MwcMisOlder--\MwcMisYounger{} Ma, including an extreme terminal stadial \citep{margari2023}. These observations establish environmental disruption within the later part of the published stasis phase.

Margari and colleagues infer that the cooling likely caused European hominin depopulation, combining proxy evidence, habitat modeling and the occupation record \citep{margari2023}. The demographic consequence is an inference with regional scope: European absence or recolonization does not establish extinction across Africa and Asia. Nor does an occupation gap directly measure the population contributing to later lineages. Environmental change could alter which populations survive, move or enter the fossil sample, complicating an anatomical trajectory even without a uniform change within each lineage.

A separate genetic analysis inferred an ancestral bottleneck at approximately \MwcBottleneckOlder--\MwcBottleneckYounger{} Ma, with effective population size near \MwcEffectiveSize{} \citep{hu2023}. Effective size is a genetic-model parameter, not a census of surviving humans. Deng and colleagues argue that the severe bottleneck can be a statistical artifact; Cousins and Durvasula find insufficient evidence and a better fit from a model without it \citep{deng2025,cousins2025}. Zhou and colleagues' July 2026 defense challenges computational and modeling criticisms, but does not establish consensus \citep{zhou2026}. Muttoni and Kent's distinct March 2024 timing-and-migration paper concerns MIS 22 near 0.9 Ma; it is neither the 2026 defense nor the MIS 34 cooling study \citep{muttoni2024}.

Australasian tektites establish an impact near \MwcImpactAge{} Ma \citep{sieh2023}. Arguments locating its source under the Bolaven volcanic field require separate geological assessment; published criticism questions that location \citep{mizera2022}. No verified evidence here establishes that the impact caused cranial-volume stasis. The impact, cooling intervals and proposed bottleneck all postdate the published phase's onset. They may contextualize its later portion or subsequent transitions, but overlap cannot explain the entire interval or identify a cause.

Mechanisms also make competing predictions. Resource stress could constrain costly development and favor energy-saving responses; environmental challenges could favor behavioral flexibility without requiring larger cranial volume. Migration, extinction and changing sampled lineages could shift the aggregate trajectory. Will and colleagues found a strong temperature association with body size, while brain-size associations were smaller and involved other environmental variables, limiting a simple cooling-to-larger-brain claim \citep{will2021}. Distinguishing these hypotheses requires geographically matched chronology, developmental and resource evidence, and within-lineage comparisons robust to sampling and dating assumptions. Establishing partner assessment's contribution would additionally require evidence for its own consequential links. The present alignment estimates none of those contributions.

\mwcfigure{minds-environment.pdf}{fig:minds-environment}{\textbf{Brain-volume trajectories and environmental context.} A: \MwcFossilZoomN{} of \MwcFossilN{} consensus observations in the enlarged window; gray points are specimens and black symbols are 100-kyr geometric means. Empty bins remain gaps. Shading marks Gingerich's published phase, not a fitted plateau; the inset shows the complete trajectory. The workbook supplies no specimen dating errors. B: the LR04 stack on its published sampling and age model; benthic $\delta^{18}$O (\textperthousand) reflects ice volume and deep-ocean temperature \citep{lisiecki2005}. C: solid bands show approximate climate-interval durations, the diamond an impact date, and the dashed hatched bar a disputed, model-inferred demographic duration. Neither band endpoints nor the bottleneck duration are confidence intervals; analytical dating uncertainties are tabulated separately in the supplement. MIS 34 cooling is proxy evidence; European hominin depopulation is Margari and colleagues' regional inference \citep{margari2023}. Alignment supplies chronological context rather than a causal estimate.}{environment}


Small-brained late-surviving hominins and contested recent brain-volume changes discourage a single anatomical ladder. The supplement retains these constraints alongside MHC preferences, genetic-lineage ratios, ADHD-associated variants and ancient games. Each supports only its specific observation, without identifying selection for cognitive ability.

The assessment account is thus an additional proposed channel within ecological, social, and cultural explanations, rather than their replacement \citep{dunbar1998,cultural,gonzalezforero2018}. It predicts that independently informative observation, the target of evaluation, and the authority governing consequential decisions should matter. Evidence for those predictions would strengthen the channel. It would still leave the magnitude of its contribution to historical cognitive and anatomical evolution to be established.

\section{Discussion and discriminating tests}\label{sec:minds-discussion}

\subsection{What the evidence supports}

The synthesis supports six distinctions: shared priorities and average trade-offs; received ratings and exposure; appearance and encounter decisions; commitment and childbearing intentions; relationship access and family size; and information gained from observers versus contexts. These results remain conditional on their measures and populations. Shared halo and institutional authority can persist as observations accumulate.

Null adjusted contrasts and country reversals narrow a claim's domain. Changes with budgeting, vote weighting or geographic recovery identify recognition, attention and selection as further questions. These contrary results remain in the figures and complete tables.

The source populations remain unlinked; conjunction of their associations cannot identify an ancestral pathway. More profiles do not repair coverage limits.

Image prediction concerns a selected, repeatedly inspected benchmark. Input tracks, rank association and score-regression $R^2$ remain distinct. Missing global duplicate components and matching prediction vectors limit the audit and any forecast of dependable contribution.

\subsection{A prospective program that can fail}

Five tests can distinguish the proposed channel from its rivals; detailed protocols appear in the supplement. First, randomize matched observation time across repeated presentation in one domain, diverse contexts, independent observers and endorsement. Test prediction of preregistered later contribution beyond baseline information. Increased liking or confidence without improved prediction would weaken the reveal mechanism. Criteria must be measured independently after the evaluation, with observation cost held fixed.

Second, separate recipient benefit from reputation. Randomize opportunities for useful work where feasible and measure delivery, judgments and subsequent access separately, alongside resources, health, training and prior ties. This distinguishes producing an immediate benefit from revealing a stable quality or receiving an institutional advantage. Third, test directed investment and compatibility in repeated dyadic interactions. Compare pair-specific forecasts with each person's average desirability, using independent delivery records to challenge a shared-halo explanation.

Fourth, measure authority: whose evaluation is used, who controls access, who can refuse, and the cost of leaving. Compare information changes with changes in access rules; consultation, consent, endorsement and coercion require separate measures. A marriage outcome can reflect practical support or restricted options without more accurate assessment. Fifth, follow later contact, reciprocal partnership, relationship quality, parenthood, births and child welfare as distinct outcomes. Retain nonparticipants and unsuccessful pairings, account for attrition, and allow appropriate follow-up periods. Restricting analysis to successful couples would condition on the transition under investigation. Each test can fail while other parts of the proposed chain remain plausible.

Historical attribution requires consequential assessment in the relevant populations, inherited variation, and a response compatible with developmental and fossil constraints. Modern associations and mathematical possibilities do not identify those facts. Ecological, demographic, cultural and institutional explanations need discriminating predictions.

\subsection{A focused connection to AI evaluation}

The AI connection concerns evaluation design: shared judge error and recognizable tests can limit forecasts of dependable benefit \citep{gao2023,sharma2024}. Evaluation should name recipient, outcome, context and authority. Broader task suites can reward common capability under Eq.~\eqref{eq:minds-breadth} while sharing one rubric. Agreement requires independent validation. These analogies do not pool or cross-validate the two papers' empirical evidence.

Choosing a future partner makes the timing of evidence especially clear. The empirical commitment is to keep assessment stages visible; the theoretical commitment is to distinguish information, contribution, and opportunity. Historical interpretation must then test every further link to inherited change.

\section*{Data, code, and disclosure}

Hash manifests preserve supplied originals. Revised human sources are under \texttt{docs/sections/brain/}, and validated aggregate evidence under \texttt{evidence/brain/}; the supplement provides provenance, rerun commands and limitations. Platform identifiers, coordinates and images remain at source. These exploratory secondary analyses recruit no new participants. Concept and direction follow the supplied Taylor--Beilharz manuscripts. Synthesis, mathematics, code and drafting were AI-assisted and require author review. Text-only author information and the full disclosure appear in the supplement; no additional affiliations or declarations are inferred.
\label{minds-main-last}
